\section{Introduction}

Automated vulnerability detection has become a proving ground for graph
learning: a function is parsed into a graph, a network is trained over it, and
detection accuracy is reported. Progress is usually claimed architecturally ---
a different message-passing scheme, an attention variant, a language-model
hybrid. Yet replication studies keep reaching the same conclusion: whatever the
architecture, performance collapses on projects the model has not
seen~\cite{chakraborty2022deep, chakraborty2024realvul, ding2025primevul}.

This paper examines a factor that most of that work holds constant: what a
node in the graph actually \emph{carries}. Nearly all graph-based detectors
embed the node's lexical content --- an identifier, an API name, a literal.
That vocabulary is a property of a codebase, not of a language. FFmpeg
allocates with \texttt{av\_malloc}, GLib with \texttt{g\_malloc}, the Linux
kernel with \texttt{kmalloc}. A model whose features come from one project's
vocabulary faces largely out-of-vocabulary input on another. The mechanism it
must learn --- an allocation whose size is computed, an indexed write, a
missing bound --- is shared; the tokens denoting it are not.

We therefore compare node representations under strict control, holding graph
topology, backbone, depth, hidden width, readout, optimiser and gradient
budget fixed and varying only the node features. Two project-invariant
alternatives to lexical tokens are considered:

\begin{itemize}
    \item an \textbf{engineered operation taxonomy} $\phi$, mapping nodes onto
          eight operation types via API allow-lists, identifier-component
          rules and syntax-directed cases --- the kind of abstraction such work
          typically constructs; and
    \item the \textbf{parser's own grammar node kinds} $\kappa$, an alphabet
          fixed by the C grammar that requires no design effort whatsoever.
\end{itemize}

The result was not the one we set out to demonstrate. Grammar kinds improve
cross-project ranking substantially and consistently: $+0.078$ AUC
($p = 0.007$) and $+0.071$ AUPRC ($p = 0.002$) over lexical features, winning
every project-level fold. The engineered taxonomy is statistically
indistinguishable from lexical features ($p \approx 0.9$), and combining it
with grammar kinds \emph{reduces} precision relative to grammar kinds alone.
For scale, the three message-passing schemes we test differ by less than half
as much under a fixed representation as the representations differ under a
fixed architecture.

We then explain the null result rather than merely reporting it. Measuring the
information the taxonomy carries relative to the grammar alphabet shows
$H(\phi \mid \kappa) = 0.06$ bits: $\phi$ is predictable from the node kind
alone for $98.2\%$ of nodes, because most of its rules are syntax-directed and
therefore already determined by the parser. Its residual contribution is
confined to two places --- which callee a call expression names, and which
operator a binary expression uses. An abstraction that is $96\%$ redundant
with a free alternative cannot be expected to add much, and does not.

Separately, we evaluate $\phi$ \emph{as a static analysis}, annotating 212
call sites and reporting per-class precision and recall. Abstractions of this
kind are routinely introduced as preprocessing and never validated. That
exercise was not ceremonial: it found that substring-based recognition of
wrapped allocators labels \texttt{\_\_put\_user} and \texttt{skb\_put}, both
writes, as release operations; that classifying \texttt{printf} and
\texttt{fprintf} as buffer-writing formatters counts 655 stream-I/O sites in
one corpus as memory operations; and that a permissive copy rule captures SIMD
lane extraction. All three are invisible to downstream accuracy.

Our contributions are:

\begin{itemize}
    \item \textbf{A controlled comparison of node representations}, with every
          other factor held fixed, across two corpora and three
          message-passing schemes (Sections~\ref{sec:rq_representation}
          and~\ref{sec:rq_generality}).
    \item \textbf{The finding that grammar node kinds suffice}, outperforming
          both lexical features and an engineered semantic taxonomy, at zero
          design cost.
    \item \textbf{An information-theoretic explanation} of why the engineered
          taxonomy adds nothing, and a diagnostic other work can reuse before
          crediting an abstraction with an effect (Section~\ref{sec:rq_why}).
    \item \textbf{Validation of the abstraction as an analysis}, with per-class
          precision and recall and the defects it uncovered
          (Section~\ref{sec:rq_analysis}).
\end{itemize}

We report the negative result at the same length as the positive one, because
the engineered taxonomy is what we built first and expected to work, and
because the redundancy diagnostic that explains its failure is cheap enough
that any similar paper could run it.
